\chapter{Implementation}
\label{chap:impl}

\section{Navigation}

\section{Development of the KPI Streaming Service}
\subsection{Technology Selection and Justification}
Technology choice followed a focused comparison based on five criteria: development speed, operational simplicity, asynchronous support, query flexibility, and suitability for small teams. 
In practical terms, the stack had to cover event ingestion, stream transport, minute and hour aggregation, and responsive API delivery for dashboard clients.

\begin{longtable}{|p{2.0cm}|p{3.5cm}|p{4.0cm}|p{4.0cm}|}
\caption[Comparison of Python API Frameworks]{Comparison of Python API Frameworks for Ingestion and Serving}
\label{table:framework_comparison}\\
\hline
\textbf{Option} & \textbf{Strengths} & \textbf{Limitations} & \textbf{Decision Note} \\
\hline
\endfirsthead
\multicolumn{4}{@{}l}{} \\
\hline
\textbf{Option} & \textbf{Strengths} & \textbf{Limitations} & \textbf{Decision Note} \\
\hline
\endhead
FastAPI & Async support, typed schemas, built in OpenAPI & Smaller ecosystem than Django & Selected for clear contracts and low integration effort \\
\hline
Flask & Very lightweight, simple routing & More manual work for validation and docs & Not selected due to higher implementation overhead \\
\hline
Django & Rich ecosystem, admin tools & Heavier stack for this scope & Not selected because project did not require full web framework features \\
\hline
\end{longtable}

This comparison showed that framework choice affected not only coding speed, but also how quickly API contracts could be documented and shared with frontend developers.

\begin{longtable}{|p{2.6cm}|p{3.4cm}|p{3.6cm}|p{4.0cm}|}
\caption[Comparison of Message Broker Alternatives]{Comparison of Message Broker Alternatives}
\label{table:broker_comparison}\\
\hline
\textbf{Option} & \textbf{Strengths} & \textbf{Limitations} & \textbf{Decision Note} \\
\hline
\endfirsthead
\multicolumn{4}{@{}l}{} \\
\hline
\textbf{Option} & \textbf{Strengths} & \textbf{Limitations} & \textbf{Decision Note} \\
\hline
\endhead
Kafka (KRaft) & High throughput, durable log, clear topic model & Requires more memory than lightweight brokers & Selected for stream oriented processing and robust replay model \\
\hline
RabbitMQ & Simple queue patterns, broad adoption & Less natural fit for log style analytics pipelines & Not selected due to weaker fit for event stream replay \\
\hline
NATS & Lightweight and fast, easy setup & Fewer built in features for long retention analytics & Not selected because durability model was less aligned with requirements \\
\hline
\end{longtable}

The comparison in Table~\ref{table:broker_comparison} shows that the broker was selected according to the requirements 
of analytical stream processing rather than throughput alone. Since the system had to retain events, support replay, 
and preserve a clear stream structure for KPI aggregation, Kafka offered the most appropriate combination of durability and processing model. 
RabbitMQ and NATS remain effective solutions for other integration scenarios, 
but in this project they were less suitable for reliable historical processing and repeatable analysis.

After selecting the transport layer, storage technology was evaluated with focus on SQL support, query flexibility, and operational overhead for a small deployment.

\begin{longtable}{|p{2.6cm}|p{3.3cm}|p{3.7cm}|p{4.0cm}|}
\caption[Comparison of Storage Alternatives]{Comparison of Storage Alternatives for KPI Facts and Aggregates}
\label{table:storage_comparison}\\
\hline
\textbf{Option} & \textbf{Strengths} & \textbf{Limitations} & \textbf{Decision Note} \\
\hline
\endfirsthead
\multicolumn{4}{@{}l}{} \\
\hline
\textbf{Option} & \textbf{Strengths} & \textbf{Limitations} & \textbf{Decision Note} \\
\hline
\endhead
PostgreSQL + TimescaleDB & SQL compatibility, time series extensions, good tooling & Performance tuning needed for heavy workloads & Selected for balance of flexibility, overall performance, and popularity \\
\hline
InfluxDB & Native time series focus, compact storage & Different query language and migration effort & Not selected to avoid SQL mismatch, migration problems, and extra learning cost for integration teams \\
\hline
ClickHouse & Fast analytical queries at scale & Higher operational complexity for SMBs & Not selected because application targeted small team operation \\
\hline
\end{longtable}

The comparison in Table~\ref{table:storage_comparison} shows that the storage layer was chosen with emphasis on query flexibility, 
integration convenience, and suitability for time-based analytical workloads. 
PostgreSQL with TimescaleDB provided the best balance between standard SQL support, time-series extensions, 
and familiar tooling, which was important for both implementation and further maintenance of the system. 
InfluxDB and ClickHouse offered strong advantages in specialized scenarios, 
but for this project they were less appropriate because of higher adaptation costs or operational complexity.

The final implementation decision also depended on deployment complexity, because the target team had limited operations capacity.

\begin{longtable}{|p{2.8cm}|p{3.4cm}|p{3.6cm}|p{3.8cm}|}
\caption[Deployment and Configuration Strategy Comparison]{Deployment and Configuration Strategy Comparison}
\label{table:deployment_strategy_comparison}\\
\hline
\textbf{Option} & \textbf{Strengths} & \textbf{Limitations} & \textbf{Decision Note} \\
\hline
\endfirsthead
\multicolumn{4}{@{}l}{} \\
\hline
\textbf{Option} & \textbf{Strengths} & \textbf{Limitations} & \textbf{Decision Note} \\
\hline
\endhead
Docker Compose + Dynaconf TOML & Fast onboarding, readable config, easy local reproduction & Limited orchestration features for large clusters & Selected as best fit for SMB setup and thesis prototype scope \\
\hline
Kubernetes + Helm & Strong scaling and control features & High operational complexity for small teams & Not selected due to significant setup and maintenance cost \\
\hline
Env only configuration & Minimal files and direct runtime mapping & Lower readability for larger parameter sets & Not selected since TOML improved clarity for many parameters \\
\hline
\end{longtable}

The comparison in Table~\ref{table:deployment_strategy_comparison} shows that deployment strategy was evaluated primarily in terms of setup effort, 
configuration clarity, and suitability for a small operational team. 
Docker Compose combined with Dynaconf TOML provided the most practical balance between reproducibility, readability, 
and ease of local and test environment management. Kubernetes with Helm and env-only configuration remain valid approaches in other contexts, 
but for this project they introduced either excessive operational overhead or lower transparency of configuration.

The comparison results led to a stack composed of FastAPI, Kafka in KRaft mode, TimescaleDB, Docker Compose, and Dynaconf based configuration. This set of tools provided a balanced compromise between implementation speed and runtime stability for the target scenario.

\subsection{Backend Service Implementation}
Implementation followed a linear backend flow from ingestion to KPI delivery. At the entry point, the ingest service accepted order and session events, validated payloads, enforced API key checks, and published normalized messages into Kafka topics. Orders were routed to \texttt{orders}, while session steps were routed to \texttt{sessions}.

Downstream processing was handled asynchronously by the stream processor. Core operations included UTC time normalization, short window deduplication by event identifier, and periodic aggregate flushes. Since aggregation relied on event time, the resulting minute and hour buckets remained consistent even when events arrived with moderate delay.

Data persistence was split into fact and aggregate relations. Raw records were stored in \texttt{orders} and \texttt{sessions}; aggregate values were stored in \texttt{kpi\_minute} and \texttt{kpi\_hour}. Derived indicators such as conversion rate were exposed through SQL views, which avoided redundant storage.

Alert generation ran in an independent service loop. It loaded recent buckets, computed smoothed current values, compared those values against a seasonal baseline, and inserted alerts only when persistence and threshold conditions were satisfied. With the current settings, both positive and negative deviations were considered once absolute deviation reached the configured boundary.

Repeatable traffic for testing came from the simulator service. It produced realistic order and session patterns, switched between normal, peak, and quiet profiles, and injected scheduled anomalies when required. This setup enabled consistent comparative experiments.

\begin{longtable}{|p{3.0cm}|p{4.9cm}|p{2.6cm}|p{2.9cm}|}
\caption[KPI Definitions Implemented in the Backend]{KPI Definitions Implemented in the Backend}
\label{table:kpi_definitions}\\
\hline
\textbf{KPI} & \textbf{Formula / Logic} & \textbf{Primary Source} & \textbf{Granularity} \\
\hline
\endfirsthead
\multicolumn{4}{@{}l}{} \\
\hline
\textbf{KPI} & \textbf{Formula / Logic} & \textbf{Primary Source} & \textbf{Granularity} \\
\hline
\endhead
Revenue & Sum of \texttt{amount} over order events & \texttt{orders} & Minute, Hour \\
\hline
Order Count & Number of order records & \texttt{orders} & Minute, Hour \\
\hline
Session Count & Number of session records with \texttt{event\_type = view} & \texttt{sessions} & Minute, Hour \\
\hline
Checkout Count & Number of session records with \texttt{event\_type = checkout} & \texttt{sessions} & Minute, Hour \\
\hline
Purchase Count & Number of session records with \texttt{event\_type = purchase} & \texttt{sessions} & Minute, Hour \\
\hline
Conversion Rate & \texttt{purchase\_count / session\_count} with zero guard & KPI Views & Minute, Hour \\
\hline
\end{longtable}

Table 4.5 summarizes KPI definitions used in aggregation logic and clarifies how each indicator was derived from fact tables or views. 
The table also demonstrates that the selected KPI set combines direct operational metrics, such as revenue and order count, 
with behavioral indicators derived from session activity. 
This structure is important because it allows the system to capture both transaction outcomes and user interaction dynamics within the same analytical pipeline. 
To support this table, it is useful to include a visual flow of backend components.

\begin{center}
\vspace{0.4em}
Recommended content: a service interaction diagram showing ingest API, Kafka topics, stream processor, TimescaleDB storage, and alerting loop.

Fig. 4.1. Backend service interaction flow
\end{center}

\subsection{API for Frontend Consumption}
The API layer exposed a compact endpoint set for dashboard use. It covered event ingestion, KPI series retrieval, latest point access, alert retrieval, and operational metrics such as freshness and time to signal. Every endpoint required API key authentication, while Swagger documentation included explicit argument and response descriptions to simplify frontend integration.

For KPI retrieval, optional segmentation by \texttt{channel} and \texttt{campaign} was supported. Time windows were accepted in ISO UTC format. If range parameters were omitted, default intervals were applied to keep common dashboard queries concise.

\begin{longtable}{|p{3.9cm}|p{3.2cm}|p{4.7cm}|p{2.3cm}|}
\caption[Backend API Endpoints]{Backend API Endpoints Exposed for Frontend Integration}
\label{table:backend_api_endpoints}\\
\hline
\textbf{Endpoint} & \textbf{Method} & \textbf{Purpose} & \textbf{Auth} \\
\hline
\endfirsthead
\multicolumn{4}{@{}l}{} \\
\hline
\textbf{Endpoint} & \textbf{Method} & \textbf{Purpose} & \textbf{Auth} \\
\hline
\endhead
\texttt{/events/order} & POST & Ingest order event into stream & API Key \\
\hline
\texttt{/events/session} & POST & Ingest session step into stream & API Key \\
\hline
\texttt{/kpi/latest} & GET & Return latest KPI point for minute or hour bucket & API Key \\
\hline
\texttt{/kpi/minute} & GET & Return minute KPI series with optional filters & API Key \\
\hline
\texttt{/kpi/hour} & GET & Return hour KPI series with optional filters & API Key \\
\hline
\texttt{/alerts} & GET & Return alerts for selected interval & API Key \\
\hline
\texttt{/metrics/freshness} & GET & Return freshness indicators per stream & API Key \\
\hline
\texttt{/metrics/time-to-signal} & GET & Return average and max signal delay & API Key \\
\hline
\end{longtable}

The endpoint matrix above can be complemented with one documentation screenshot to show how parameters and response contracts are presented to frontend developers in practice.

\begin{center}
\vspace{0.4em}
Recommended content: API documentation page with endpoint summary, query arguments, and response model examples.

Fig. 4.3. API documentation view for frontend integration
\end{center}

\subsection{Deployment and Runtime Configuration}
Deployment at runtime was orchestrated with Docker Compose. The service set included Kafka, topic initialization, TimescaleDB, ingest API, stream processor, alerting, and simulator components. Health checks and dependency conditions reduced startup race risks by delaying dependent services until core components became healthy.

Configuration management separated defaults from local secrets. Baseline values were stored in \texttt{settings.toml}, whereas sensitive and environment specific overrides were placed in \texttt{.secrets.toml}. This arrangement improved readability and made deployments reproducible across environments.

During operational tests on small virtual machines, memory pressure occasionally terminated Kafka. As a mitigation, runtime guidance included memory monitoring, lighter simulator profiles, and swap activation for low memory hosts.

The most relevant runtime parameters were grouped by service. For the ingest API, \texttt{API\_KEY} controlled request authorization and protected endpoints used by simulator and frontend clients. In the stream processor, \texttt{FLUSH\_INTERVAL\_\allowbreak SECONDS} defined how frequently aggregates were flushed to database tables, which directly affected dashboard update cadence. In alerting, \texttt{THRESHOLD\_PCT} and \texttt{DURATION\_\allowbreak MINUTES} controlled sensitivity and persistence, reducing noisy signals by requiring sustained deviation. In the simulator, \texttt{SEND\_INTERVAL\_\allowbreak SECONDS} regulated traffic pace, while \texttt{SCHEDULE\_MODE} switched between workload patterns such as \texttt{day-night} and \texttt{seasonal}. Together, these parameters provided practical control over stability, responsiveness, and experimental repeatability.



\subsection{Dashboard and Interaction Design for Real-Time KPI Monitoring}
% Design-focused subsection written by a separate contributor.

\subsection{Frontend Application Implementation}
% Frontend-focused subsection written by a separate contributor.

\section{Data Collection}
% Describe how runtime and KPI-related data were collected during experiments.

\section{Data Validation}
% Describe validation workflow: expected values, deduplication checks,
% segment-level checks, anomaly checks, and consistency checks across runs.

\section{Summary}

